\section*{Hoofdstuk \ref{ch:b2:weather-systems} --- Weersystemen}

\begin{solution}{exo:b2:weather-systems:1}
(a) $f = 2 \times 7.292 \times 10^{-5} \times \sin\ang{50} = \qty{1.117e-4}{s^{-1}}$. (b) Met
$\bar T = \qty{278}{K}$ is $\tan\alpha = 1.117 \times 10^{-4} \times 278 \times 15/(9.81 \times 6)
= 0.00791$, ongeveer 1 op 126. (c) $1/0.00791 = 126$: zo'n \qty{126}{km} achter het front. (d) De
helling is omgekeerd evenredig met het temperatuurcontrast: ze halveert, tot 1 op 253.
\end{solution}

\begin{solution}{exo:b2:weather-systems:2}
(a) $6 \times 150 = \qty{900}{km}$ ervoor. (b) $900/25 = \qty{36}{h}$, anderhalve dag. (c) Eerst
de hoge cirrus, dan cirrostratus, dan altostratus; de bewolking wordt dikker en lager, en regen
valt uit de lage, dikke bewolking terwijl het front nadert.
\end{solution}

\begin{solution}{exo:b2:weather-systems:3}
De symbolen wijzen in de richting waarin het front beweegt, en in de warmere lucht vóór een koufront of de
koudere lucht vóór een warmtefront. (a) Een koufront: de koude lucht ligt ten noordwesten ervan, erachter;
het beweegt naar het zuidoosten. (b) Een warmtefront: de koude lucht ligt ten noorden ervan, ervoor; het beweegt
naar het noorden. (c) Een \omterm{def:b2:weather-systems:fronts}{occlusiefront}, dat het lagedrukgebied verbindt met het tripelpunt waar het koufront en het
warmtefront beginnen; koude lucht ligt aan beide kanten ervan, de warme lucht is omhooggetild; het beweegt in de richting waarin zijn
symbolen wijzen. (d) Een \omterm{def:b2:weather-systems:fronts}{stationair front}: de driehoeken wijzen naar het zuiden, in de warme lucht, dus de koude
lucht ligt ten noorden; het beweegt nauwelijks.
\end{solution}

\begin{solution}{exo:b2:weather-systems:4}
(a) $\beta = 2 \times 7.292 \times 10^{-5} \times \cos\ang{45}/(6.371 \times 10^6) =
\qty{1.62e-11}{m^{-1}.s^{-1}}$. (b) Voor \qty{4000}{km} is $k = \qty{1.571e-6}{m^{-1}}$, $\beta/k^2 =
\qty{6.6}{m/s}$ en $c = 12 - 6.6 = \qty{5.4}{m/s}$; voor \qty{8000}{km} is $\beta/k^2 =
\qty{26.3}{m/s}$ en $c = \qty{-14.3}{m/s}$. (c) De korte golf beweegt naar het oosten, \qty{470}{km} per dag;
de lange beweegt naar het westen, \qty{1240}{km} per dag, tegen de wind in.
\end{solution}

\begin{solution}{exo:b2:weather-systems:5}
(a) $c = 0$ wanneer $U = \beta/k^2$, dus $k = \sqrt{\beta/U}$ en $\lambda = 2\pi/k =
2\pi\sqrt{U/\beta}$. (b) $2\pi\sqrt{12/(1.62 \times 10^{-11})} = \qty{5410}{km}$. (c) De
breedtecirkel van \ang{45} meet $2\pi \times 6371 \times \cos\ang{45} = \qty{28300}{km}$: ongeveer 5.2
golven, dus in de praktijk vijf. (d) $\lambda$ groeit als $\sqrt U$: een sterkere westenwind maakt de
stilstaande golven langer, en er passen er minder rond het halfrond.
\end{solution}

\begin{solution}{exo:b2:weather-systems:6}
(a) In $C_+ = (\sqrt{72}, \sqrt{72}, 27)$: $\dot x = 10(\sqrt{72} - \sqrt{72}) = 0$, $\dot y =
\sqrt{72}(28 - 27) - \sqrt{72} = 0$, $\dot z = 72 - \tfrac83 \times 27 = 72 - 72 = 0$.
(b) $\mathrm{e}^{-13.67} = 1.16 \times 10^{-6}$: het volume krimpt in één tijdseenheid bijna een miljoen keer.
(c) Een volume kan krimpen terwijl het in één richting wordt uitgerekt en in de andere veel
sterker wordt samengedrukt: de afstand langs de uitgerekte richting groeit, en twee naburige
banen lopen erlangs uiteen, terwijl het volume ertussen verdwijnt.
\end{solution}

\begin{solution}{exo:b2:weather-systems:7}
(a) $\ln 2/0.906 = 0.765$ tijdseenheden. (b) $\ln(10/10^{-8})/0.906 = \ln(10^9)/0.906 = 22.9$
tijdseenheden. (c) $\ln 1000/0.906 = 7.6$ tijdseenheden: een duizendvoudige verbetering van de begintoestand
koopt maar een derde van de tijd uit (b), want $\ln 10^3/\ln 10^9 = 1/3$.
\end{solution}

\begin{solution}{exo:b2:weather-systems:8}
(a) $(303 - 195)/195 = 0.554$ en $V_{\max} = \sqrt{0.9 \times 0.554 \times 15\,000} =
\qty{86.5}{m/s}$. (b) Ruim boven de \qty{70}{m/s} van categorie 5, dicht bij de \qty{86}{m/s} die
vijf recente stormen hebben overschreden. (c) Stormen blijven zelden lang in de stationaire toestand die de theorie
aanneemt; verticale windschering, droge lucht die de storm in wordt gezogen en de dalende lucht eromheen, die
de theorie weglaat, verzwakken ze.
\end{solution}

\begin{solution}{exo:b2:weather-systems:9}
(a) $(303 - 195)/303 = 0.356$. (b) $(303 - 195)/195 = 0.554$; de verhouding is $0.554/0.356 =
303/195 = 1.554$. (c) In een carnotmotor verlaat de arbeid de motor. In een tropische cycloon wordt de
arbeid door wrijving dicht bij het zeeoppervlak gedissipeerd, en die warmte verwarmt de instromende lucht bij de
temperatuur $T_s$, waar de motor zijn warmte opneemt: ze wordt opnieuw gebruikt, en de motor gedraagt zich
alsof hij efficiënter was met de factor $T_s/T_o$.
\end{solution}

\begin{solution}{exo:b2:weather-systems:10}
(a) De afleiding van \cref{prop:b2:weather-systems:slope} zonder benadering geeft
$\mathrm{d}z/\mathrm{d}x = f(\rho_c v_c - \rho_w v_w)/[g(\rho_c - \rho_w)]$. Een dichtere vloeistof
die op een lichtere rust, is onstabiel: de kleinste verstoring doet haar omslaan. In een front in
rust moet de koude, dichtere lucht daarom onder liggen, wat de wig van de definitie is; dan is
$\rho_c > \rho_w$, stijgt het vlak naar de koude lucht toe, en is $\rho_c v_c > \rho_w v_w$.
(b) $\zeta = \partial v/\partial x = 10/(50 \times 10^3) = \qty{2e-4}{s^{-1}}$, ongeveer $1.9f$ met
$f = \qty{1.03e-4}{s^{-1}}$. (c) Positief op het noordelijk halfrond: cyclonaal, en even groot als
de Coriolisparameter zelf; fronten zijn banden van sterke cyclonale vorticiteit.
\end{solution}

\begin{solution}{exo:b2:weather-systems:11}
(a) $\log_2(6/0.5) = \log_2 12 = 3.58$ verdubbelingen, $3.58 \times 2.4 = 8.6$~dagen. (b) Eén
dag verwachting meer vraagt één verdubbelingstijd minder groei voordat de fout onbruikbaar is: de
beginfout moet elk decennium met $2^{1/2.4} = 1.33$ krimpen. (c) Een vertienvoudigde verkleining is
$\log_2 10 = 3.32$ verdubbelingen; als die maar 4 tot 5 dagen wint, verdubbelen kleine fouten in $4/3.32 = 1.2$
tot $5/3.32 = 1.5$ dagen: kleinere fouten groeien sneller dan de grotere die in de
verwachtingen van de jaren 1980 werden gemeten, door latente warmte die op kleine schaal in convectie vrijkomt.
\end{solution}

\begin{solution}{exo:b2:weather-systems:12}
(a) Bij \qty{30}{\celsius} is $e_s = \qty{43.5}{hPa}$ en $q_s^* = \qty{27.1}{g/kg}$; de lucht van
\qty{29}{\celsius} heeft $e = 0.8 \times 41.0 = \qty{32.8}{hPa}$ en $q_a = \qty{20.4}{g/kg}$;
$h_o^* - h^* = 1004.5 \times 1 + 2.501 \times 10^6 \times 0.0067 = \qty{17.8}{kJ/kg}$, tegen
\qty{16.8}{kJ/kg}. (b) $(303.15 - 200)/200 = 0.516$ en $V_{\max} = \sqrt{0.9 \times 0.516
\times 17\,800} = \qty{90.9}{m/s}$, tegen \qty{87.9}{m/s}. (c) Ongeveer \qty{3.4}{\percent} per
kelvin. Het onevenwicht stijgt met \qty{6}{\percent}, wat \qty{3}{\percent} aan de snelheid toevoegt;
de factor $(T_s - T_o)/T_o$ stijgt met \qty{1}{\percent}, wat \qty{0.5}{\percent} toevoegt: de
brandstof, het vocht dat de warmere zee kan leveren, draagt het meest bij, via de exponentiële stijging van
$e_s$.
\end{solution}

\begin{solution}{pb:b2:weather-systems:1}
\textbf{1.} $(480 - 390)/6 = \qty{15}{km/h}$.
\textbf{2.} $390/15 = \qty{26}{h}$ na 12 UTC: om 14 UTC de volgende dag.
\textbf{3.} $f = \qty{1.117e-4}{s^{-1}}$.
\textbf{4.} Met $\bar T = \qty{281}{K}$ is $\tan\alpha = 1.117 \times 10^{-4} \times 281 \times
16/(9.81 \times 8) = 0.0064$, ongeveer 1 op 156.
\textbf{5.} Het front ligt dan $6 \times 15 = \qty{90}{km}$ voorbij de stad; het vlak ligt
$90\,000 \times 0.0064 = \qty{576}{m}$ erboven.
\textbf{6.} $1500/0.0064 = \qty{234}{km}$ achter het front, bereikt $234/15 = \qty{15.6}{h}$ na
de passage.
\textbf{7.} Cumulus en cumulonimbus, buien, misschien onweer, windstoten en een plotselinge daling van de
temperatuur; de druk, die daalt terwijl het front nadert, begint te stijgen wanneer het passeert; daarna
koelere, helderdere lucht.
\textbf{8.} $\beta = 2 \times 7.292 \times 10^{-5} \times \cos\ang{50}/(6.371 \times 10^6) =
\qty{1.47e-11}{m^{-1}.s^{-1}}$.
\textbf{9.} $k = 2\pi/(5 \times 10^6) = \qty{1.257e-6}{m^{-1}}$ en $\beta/k^2 = \qty{9.3}{m/s}$.
\textbf{10.} $c = 20 - 9.3 = \qty{10.7}{m/s}$, oostwaarts.
\textbf{11.} $10.7 \times 86\,400 = \qty{924}{km}$ naar het oosten.
\textbf{12.} $\lambda_s = 2\pi\sqrt{20/(1.47 \times 10^{-11})} = \qty{7330}{km}$; de
breedtecirkel van \ang{50} meet $2\pi \times 6371 \times \cos\ang{50} = \qty{25700}{km}$, ongeveer 3.5
zulke golven.
\textbf{13.} Ervoor, aan zijn oostkant, waar de korte golven stijgende lucht opwekken en de
druk aan het oppervlak daalt.
\textbf{14.} $\lambda = \ln(3.9 \times 10^{-3}/5 \times 10^{-7})/(14 - 4) = \ln 7800/10 = 0.896$.
\textbf{15.} $\ln 2/0.896 = 0.77$ tijdseenheden.
\textbf{16.} \qty{1}{\percent} onder 0.906: het gemiddelde van 200 paren, een eindige tijd en aflezingen uit
een grafiek.
\textbf{17.} Zodra de twee runs ver uit elkaar liggen, bevinden ze zich op onafhankelijke punten van de attractor, en
kan hun afstand zijn grootte niet overschrijden: de groei verzadigt.
\textbf{18.} $\ln 10/0.896 = 2.6$ tijdseenheden later.
\textbf{19.} $\log_2(3/0.1) = \log_2 30 = 4.91$ verdubbelingen.
\textbf{20.} $4.91 \times 2.4 = 11.8$~dagen.
\textbf{21.} Eén verdubbelingstijd: 2.4~dagen.
\textbf{22.} $4.91 \times 5 = 24.5$~dagen: de vroege schatting was veel te optimistisch.
\textbf{23.} Met een verdubbelingstijd van 2.4 dagen zou een tienmaal kleinere fout $3.32 \times 2.4
= 8.0$~dagen winnen, niet 4 tot 5: kleine fouten groeien sneller dan 2.4 dagen per verdubbeling, en een
schatting die op die verdubbelingstijd is gebouwd, overschat de horizon.
\textbf{24.} Deze verwachting zou hoogstens ongeveer \textbf{12 dagen} bruikbaar kunnen zijn, en in de praktijk
korter.
\end{solution}
